\documentclass{article}
\usepackage{hyperref}
\usepackage{amsmath}
\usepackage{graphicx}

\title{Sample DSR Paper}
\author{Playground CLI}
\date{\today}

\begin{document}
\maketitle

\begin{abstract}
This DSR paper demonstrates the clap-noun-verb framework for building semantic CLIs.
\end{abstract}

% Design Science Research (DSR) Thesis Structure: Problem, Artifact, Evaluation, Theory
% This structure follows the DSR methodology for creating and evaluating artifacts

\section{Problem}
Identification of the research gap and motivation.

Identification and motivation of the research problem, establishing relevance and defining the problem space for design science inquiry.

\section{Artifact}
Design and implementation of the solution artifact.

Design and development of the solution artifact, including design principles, architecture, and implementation details.

\section{Evaluation}
Validation of artifact effectiveness.

Rigorous evaluation of the artifact's effectiveness and utility, demonstrating that it addresses the identified problem.

\section{Theory}
Contribution to the knowledge base.

Contribution to the theoretical knowledge base, generalizing findings and establishing design principles for future research.

\bibliographystyle{plain}
\bibliography{references}

\end{document}
